\section{User Interfaces and Supported Use Cases}
\label{sec:iaui-from-prototypes}

The prototypes developed in the previously published works do not only validate the architectural assumptions, but also demonstrate how the \ac{dia} can concretely support real user-facing tasks. In this section, we highlight the \acp{ui} that have been actually implemented in the three prototypes and we position them with respect to the five use cases introduced in~\Cref{sec:intro}.

\paragraph{Visualization and inspection interfaces.}
All three prototypes~\cite{BELLANDI2024,pozzi_wise_2025,Agazzi2025DAVE} include a \emph{document explorer} to visualize annotated documents and inspect entity mentions.
These interfaces support \ref{uc:navigation} and, to some extent, enable manual verification and minor correction actions directly from the viewer.

\paragraph{Search and investigative retrieval.}
\textcite{pozzi_wise_2025} introduced a faceted-search interface for investigative retrieval over IMA data, covering \ref{uc:investigative-retrieval}.
A conceptually analogous faceted search was later adopted in~\cite{Agazzi2025DAVE} for \ref{uc:case-retrieval}, demonstrating the reuse of the same interaction paradigm across domains.

\paragraph{Graph-based exploration.}
The WISE\todo{remove wise} prototype~\cite{pozzi_wise_2025} also included a graph-based interface tailored to investigative analysis of communication networks derived from chats, supporting complex instances of \ref{uc:investigative-retrieval} that cannot be expressed through keyword search alone.

\paragraph{Question answering and conversational interaction.}
\textcite{Agazzi2025DAVE} introduced a conversational \ac{qa} interface supporting \ref{uc:qa}, allowing users to query document collections in natural language with linked evidence.

\paragraph{Error correction and refinement.}
All three works included partial support for human correction through the document explorer. In~\cite{Agazzi2025DAVE}, this capability was extended with an interface for refining entity clusters, moving toward structured \ref{uc:investigative-retrieval} and \ref{uc:qa} pipelines with human oversight.

\paragraph{Analytical inspection.}
In \textcite{BELLANDI2024}, statistical analysis was explored in notebooks rather than via a dedicated \ac{ui}, partially addressing \ref{uc:stats} but without a deployed interface.

\medskip

\noindent
As a synthesis, \Cref{tab:uc:ui} will map each \ac{ui} to the supported use cases \ref{uc:case-retrieval}--\ref{uc:stats}. Notably, no published prototype includes a \emph{\ac{ui} for \ref{uc:stats} statistical monitoring} or a \emph{dedicated \ac{ui} for \ref{uc:investigative-retrieval} error correction} beyond the lightweight interactions embedded in the document explorers. These uncovered areas motivated the architectural requirements for visualization and error-correction defined in this chapter.

\begin{table}[htbp]
\centering
\label{tab:uc:ui}
\begin{threeparttable}
\caption{Published \acp{ui} and supported use cases.}
\begin{tabularx}{\textwidth}{p{4cm} X p{1.5cm}}
\toprule
\textbf{\ac{ui} type} & \textbf{Supported UC / principles} & \textbf{Used in Ref.} \\
\midrule

Document explorer &
\ref{uc:navigation} (navigation); error correction (partial) &
\cite{BELLANDI2024,pozzi_wise_2025,Agazzi2025DAVE} \\

Faceted search &
\ref{uc:case-retrieval} (case retrieval);
\ref{uc:investigative-retrieval} (investigative retrieval) &
\cite{pozzi_wise_2025,Agazzi2025DAVE} \\

Graph-based exploration &
\ref{uc:investigative-retrieval} (investigative retrieval) &
\cite{pozzi_wise_2025} \\

Conversational \ac{qa} \ac{ui} &
\ref{uc:qa} (\ac{qa} over documents/collections) &
\cite{Agazzi2025DAVE} \\

Cluster refinement \ac{ui} &
error correction (partial, entity-level refinement) &
\cite{Agazzi2025DAVE} \\

Notebook-based analysi\tnote{a} &
\ref{uc:stats} (statistical analysis, partial) &
\cite{BELLANDI2024} \\
\bottomrule
\end{tabularx}

\begin{tablenotes}
    \item[a] Implemented as exploratory analysis in Jupyter notebooks \url{https://jupyter.org/}.
\end{tablenotes}

\end{threeparttable}
\end{table}


As shown in Table~\ref{tab:iaui-published}, the \acp{ui} prototyped so far do not cover all the intended use cases of the \ac{dia}. In particular, the following gaps remain:

\begin{enumerate}[label=(\roman*)]
    \item \textbf{Lack of a deployed \ref{uc:stats} interface}: statistical analysis is currently supported only through notebook-based workflows and not through a dedicated, interactive \ac{ui}.
    \item \textbf{Partial error correction}: correction is possible only indirectly through document explorers or via entity refinement, but no dedicated error-correction \ac{ui} exists.
    \item \textbf{Limited reuse of \acp{ui} across domains}: except for the faceted search (reused in both \cite{pozzi_wise_2025} and \cite{Agazzi2025DAVE}), most implemented \acp{ui} were tailored to specific datasets or tasks rather than conceived as reusable architectural components.
\end{enumerate}

These missing capabilities directly motivated the architectural requirements for visualization and error-correction interfaces defined earlier in this chapter.


\subsection{UI Prototypes in \textcite{BELLANDI2024}}
\label{sec:ui-clsr}

\paragraph{Document explorer.}
Descrivere cosa permette di vedere, su quali dati era applicato (civil judgments),
e che supporto dava per \ref{uc:navigation} e verifica.

\paragraph{Faceted / investigative search (if any).}
Breve nota se nel lavoro era presente solo a livello concettuale o implementato.

\paragraph{Statistical inspection via notebooks.}
Collegare a \ref{uc:stats}, chiarendo che non era una \ac{ui}.

\medskip
% --- next paper ---

\subsection{UI Prototypes in \textcite{pozzi_wise_2025}}
\label{sec:ui-wise}

\paragraph{Faceted investigative retrieval.}
Descrivere come veniva usata per \ref{uc:investigative-retrieval}.

\paragraph{Graph-based exploration.}
Ruolo per analisi investigativa, differenza rispetto a faceted,
e relazione con \ref{uc:investigative-retrieval}.

\paragraph{Document explorer with partial correction.}
Richiamo a \ref{uc:navigation} e ErrCorr(partial).

\medskip
% --- next ---

\subsection{UI Prototypes in \textcite{Agazzi2025DAVE}}
\label{sec:ui-dave}

\paragraph{Faceted case retrieval.}
Riutilizzo del paradigma di WISE ma su contesto legale — mapping a \ref{uc:case-retrieval}.

\paragraph{Conversational question answering.}
Collegare a \ref{uc:qa} e differenza rispetto a faceted.

\paragraph{Document explorer.}
Come in altri paper, ma arricchito.

\paragraph{Cluster refinement.}
Menzionare come abilitazione all'ErrCorr (partial).

\medskip

\subsection{Cross-paper observations}
\label{sec:ui-cross}

% Qui due o tre bullet:
- riuso della stessa interfaccia (faceted) su due domini diversi;
- funzioni presenti in tutti (document explorer) vs. funzioni emergenti (QA/refinement);
- limiti individuati (niente UI per \ref{uc:stats}, correzione mai completa).
